\scanpage{77}
\mspage{21}

i.e. une $W$-équivalence fibre par fibre, il faut \unsure{s.|?.} que $f$ soit une $W$ équivalence fibre par fibre (et cela doit impliquer \inserted{cf \textcircled{1} p.~15} que $f$ est une $W$-équivalence)

En effet, dans le diagramme commutatif
\diagram{\begin{tikzcd}[ampersand replacement=\&]
X_s\arrow[r,"f_s"]\arrow[d]\&Y_s\arrow[d]\\
X_{/s}\arrow[r,"f_{/s}"']\&Y_{/s}
\end{tikzcd}}
les flèches verticales sont dans $W$, car \deleted{\unsure{f est|sont}} \inserted{$X,Y$} propres sur $S$ \inserted{(ou $X,Y$ dans $\underline{\mathrm{Fib}\,\mathrm{hot}}\,S$)}, donc $f_s$ est dans $W$ ssi $f_{/s}$ l'est.
% margin note written sideways
\margin{\uline{Cor. 3} Si $X\in\operatorname{Ob}\underline{\mathrm{Fib}\,\mathrm{hot}}\,S$, \unsure{$\blacksquare$} alors $\Psi(X)$ est dans $\underline{\mathrm{Fib}\,\mathrm{sc}}_{\unsure{0|o}}\,S$. \deleted{\unsure{On regarde}} Cela signifie en effet que les $X_{/s}\unsure{\text{---}|\to}X_{/t}$ sont dans $W$, or l'hyp. sur $X$ \unsure{\{} \unsure{indique|implique} que c'est des $\{W_\infty\}$\unsure{,} OK.}
% margin note written sideways
\margin{\diagram{\deleted{\begin{tikzcd}[ampersand replacement=\&,column sep=small,row sep=small]
 X_s\arrow[r]\arrow[d]\&X_t\arrow[d]\\
 X_{/s}\arrow[r]\& X_{/t}
\end{tikzcd}}}}

$\underline{\Phi}$ se définit de façon duale à $\Psi$, ce qui revient à échanger les rôles de $\sigma_X,\beta_X$ dans le diagramme page 17.

II) Définition de \deleted{\unsure{Parf}} $C:\mathrm{Cat}_{/S}\unsure{\text{---}|\to}\underline{\mathrm{Parf}}\ \deleted{\unsure{X}}S$.

Cette fois, c'est la factorisation de \deleted{\unsure{toute}} flèches $\deleted{\unsure{f}}p:X\longrightarrow S$ dans $\mathrm{Cat}$, en
\[ X\xrightarrow{j_{X/S}}\overline{X}\xrightarrow{\bar p} S, \]
où $i_X$ est une Hot-équivalence (donc dans $W$), et $\bar p$ parfait. Cela se construit formellement comme tantôt, en utilisant $\underline{\mathrm{Ch}}_\infty(\deleted{\unsure{X}}S)$ au lieu de $\underline{\mathrm{Fl}}(\deleted{\unsure{X}}S)$, et la factorisation
\[ \deleted{\unsure{X}}S\xrightarrow{\deleted{\unsure{i}}j_S}\underline{\mathrm{Ch}}_\infty(\deleted{\unsure{X}}S)\xrightarrow{p_S=(\sigma_S,\beta_S)}\deleted{\unsure{X}}S\times\deleted{\unsure{X}}S \]
de $\mathrm{diag}_S:S\longrightarrow S\times S$. La situation est plus jolie, car $p_S$ est déjà ``parfait'', ce qui \unsure{impliquera|implique} formellement que le $\bar p$ qui s'en déduit, pour factoriser $p:X\longrightarrow S$, est lui-même parfait.

% archive sheet number: 76
